\documentclass[11pt]{article}
\usepackage[a4paper,margin=1in]{geometry}
\usepackage{fontspec}
\usepackage[boldfont=false]{xeCJK}
\setCJKmainfont{FandolSong-Regular.otf}
\usepackage{amsmath}
\usepackage{amssymb}
\usepackage{array}
\usepackage{booktabs}
\usepackage{graphicx}
\usepackage{adjustbox}
\usepackage{enumitem}
\setlist[itemize]{topsep=2pt,partopsep=0pt,parsep=0pt,itemsep=2pt,leftmargin=1.4em}
\setlist[enumerate]{topsep=2pt,partopsep=0pt,parsep=0pt,itemsep=2pt,leftmargin=1.6em}
\usepackage{parskip}
\usepackage{fvextra}
\fvset{breaklines=true,breakanywhere=true,fontsize=\small}
\setlength{\parindent}{0pt}

\begin{document}

\begin{center}
  {\LARGE\bfseries Group 2}\\[0.35em]
  {\large A Level Chemistry}
\end{center}
\vspace{0.6em}

\section*{The Group 2 metals}

\textbf{Group} 族 2 holds the metals magnesium, calcium, strontium and barium. They all have two outer electrons, which they lose to form $2+$ ions. Going \textbf{down} the group, the atoms get larger and the outer electrons are easier to lose, so the metals get \textbf{more reactive} — their \textbf{reactivity} 反应活性 increases down the group.

\subsection*{Reactions of the elements}

With \textbf{oxygen}, they burn to form an oxide:

\[
2\text{Mg} + \text{O}_2 \rightarrow 2\text{MgO}
\]

With \textbf{water}, they form a hydroxide and hydrogen. The reaction gets faster down the group:

\[
\text{Ca} + 2\text{H}_2\text{O} \rightarrow \text{Ca(OH)}_2 + \text{H}_2
\]

Magnesium is slow with cold water but reacts fast with steam, giving $\text{MgO}$ and $\text{H}_2$.

With \textbf{dilute hydrochloric or sulfuric acid}, they form a salt and hydrogen:

\[
\text{Mg} + 2\text{HCl} \rightarrow \text{MgCl}_2 + \text{H}_2 \qquad \text{Mg} + \text{H}_2\text{SO}_4 \rightarrow \text{MgSO}_4 + \text{H}_2
\]

With sulfuric acid the reaction slows down the group, because the \textbf{sulfates} 硫酸盐 formed (such as $\text{BaSO}_4$) are insoluble and coat the metal.

\subsection*{Reactions of the compounds}

The \textbf{oxides} 氧化物 and \textbf{hydroxides} 氢氧化物 are basic. They react with water and with dilute acids:

\[
\text{MgO} + \text{H}_2\text{O} \rightarrow \text{Mg(OH)}_2 \qquad \text{Mg(OH)}_2 + 2\text{HCl} \rightarrow \text{MgCl}_2 + 2\text{H}_2\text{O}
\]

The \textbf{carbonates} 碳酸盐 react with dilute acids to give a salt, water and carbon dioxide:

\[
\text{CaCO}_3 + 2\text{HCl} \rightarrow \text{CaCl}_2 + \text{H}_2\text{O} + \text{CO}_2
\]

\subsection*{Thermal decomposition}

\textbf{Thermal decomposition} 热分解 means breaking a compound apart by heating it.

\begin{itemize}
\item the carbonates break into the oxide and carbon dioxide:

\end{itemize}

\[
\text{CaCO}_3 \rightarrow \text{CaO} + \text{CO}_2
\]

\begin{itemize}
\item the \textbf{nitrates} 硝酸盐 break into the oxide, brown nitrogen dioxide gas and oxygen:

\end{itemize}

\[
2\text{Ca(NO}_3)_2 \rightarrow 2\text{CaO} + 4\text{NO}_2 + \text{O}_2
\]

The \textbf{thermal stability} 热稳定性 of both the carbonates and the nitrates \textbf{increases} down the group. A larger metal ion pulls less on the carbonate or nitrate ion, so the compound is harder to break apart — it needs a higher temperature. So magnesium carbonate decomposes most easily, and barium carbonate is the hardest.

\subsection*{Trends in solubility}

\begin{center}
\adjustbox{max width=\linewidth}{%
\begin{tabular}{ll}
\toprule
\textbf{Compound} & \textbf{Trend in \textbf{solubility} 溶解度 down the group} \\
\midrule
hydroxides & \textbf{increase} (Mg(OH)$_2$ almost insoluble; Ba(OH)$_2$ soluble) \\
sulfates & \textbf{decrease} (MgSO$_4$ soluble; BaSO$_4$ insoluble) \\
\bottomrule
\end{tabular}}
\end{center}

From these trends you can predict the properties of the next element down, radium, and of its compounds.


\end{document}
